Los resultados son coherentes con la l\'ogica econ\'omica del problema: viviendas con m\'as habitaciones tienden a costar m\'as, mientras que zonas con mayor proporci\'on de poblaci\'on de bajo estatus socioecon\'omico presentan menores valores de mercado. La variable \textit{chas} muestra un efecto positivo, pero con una muestra muy desbalanceada entre los casos 0 y 1, por lo que conviene interpretar la diferencia con cautela.

El PCA result\'o \textit{útil} para reducir dimensionalidad y resumir la estructura del conjunto, aunque no reemplaza el an\'alisis inferencial de variables puntuales. Su fortaleza aqu\'i es exploratoria: permite visualizar agrupamientos y detectar qu\'e variables dominan el comportamiento conjunto. Como limitaci\'on, el modelo lineal para \textit{crim} puede verse afectado por colinealidad entre predictores, por lo que una extensi\'on natural ser\'ia validar el ajuste con subconjuntos de entrenamiento/prueba o con regularizaci\'on.

Otra limitaci\'on importante es que en la pregunta de suburbios m\'as baratos no se dispone de nombres geogr\'aficos, sino de identificadores de observaci\'on; por eso el informe debe hablar de observaciones o suburbios indexados, no de barrios con nombre propio. Aun as\'i, el an\'alisis permite responder las preguntas con evidencia suficiente y reproducible.